\documentclass[a4paper,12pt]{article}
\usepackage[utf8]{inputenc}
\usepackage[T1]{fontenc}
\usepackage[ngerman]{babel}
\usepackage{amsmath}
\usepackage{amssymb}
\usepackage{enumitem}
\usepackage{geometry}
\geometry{a4paper, margin=2.5cm}

\title{Einsendeaufgaben -- Lektion 3\\[0.5em]
\large Modul 61111: Mathematische Grundlagen}
\author{}
\date{}

\begin{document}
\maketitle

\section*{Aufgabe 3.1}

\begin{enumerate}[label=\alph*)]

    \item
    \[
        a_1(-u) + a_2(u+v+w)
        = (-a_1+a_2)\,u + a_2\,v + a_2\,w = 0
    \]
    Da $u,v,w$ linear unabhängig ist, muss $a_2 = 0$ sein.
    Außerdem muss $-a_1+a_2 = 0$ sein. Wegen $a_2 = 0$
    muss $a_1 = 0$ sein.

    $\Rightarrow$ $-u$, $u+v+w$ sind linear unabhängig.

    \item
    \[
        a_1(u+v) + a_2(v+w) + a_3(v-w)
        = a_1\,u + (-a_1+a_2+a_3)\,v + (a_2-a_3)\,w
        = 0
    \]
    Wegen linearer Unabhängigkeit muss $a_1 = 0$ sein.
    Aus $-a_1+a_2+a_3 = a_2+a_3 = 0$
    und $a_2 = a_3$ wegen $a_2-a_3 = 0$
    folgt $a_2 = a_3 = 0$.

    $\Rightarrow$ $u+v$, $v+w$, $v-w$ linear unabhängig.

    \item
    \[
        a_1(u-v) + a_2(u-w) + a_3(v-w)
        = (a_1+a_2)\,u + (-a_1+a_3)\,v + (-a_2-a_3)\,w
    \]
    Gleichungssystem aufstellen:
    \[
        \begin{pmatrix} 1 & 1 & 0 \\ -1 & 0 & 1 \\ 0 & -1 & -1 \end{pmatrix}
        \begin{pmatrix} a_1 \\ a_2 \\ a_3 \end{pmatrix} = \mathbf{0}
    \]
    L\"{o}sung durch Zeilenumformungen:
    \[
        \begin{pmatrix} 1 & 1 & 0 \\ -1 & 0 & 1 \\ 0 & -1 & -1 \end{pmatrix}
        \;\Big|\; Z_{2+1}(1)
        \longrightarrow
        \begin{pmatrix} 1 & 1 & 0 \\ 0 & 1 & 1 \\ 0 & -1 & -1 \end{pmatrix}
        \;\Big|\; Z_{3+2}(1)
        \longrightarrow
        \begin{pmatrix} 1 & 1 & 0 \\ 0 & 1 & 1 \\ 0 & 0 & 0 \end{pmatrix}
    \]
    $\Rightarrow$ 2 Pivot-Positionen

    $\Rightarrow$ linear abhängig

\end{enumerate}

\end{document}
